% sec1_3.tex — Section 1.3: Finite-Sample Properties of OLS

Having derived the OLS estimator, we now examine its finite-sample properties,
namely, the characteristics of the distribution of the estimator that are valid for any
given sample size $n$.

%% ─── Finite-Sample Distribution of b ──────────────────────────────────────────
\subsection*{Finite-Sample Distribution of \texorpdfstring{$\vec{b}$}{b}}

\begin{proposition}[finite-sample properties of the OLS estimator of $\bm{\beta}$]
\label{prop:1.1}
\begin{enumerate}[label=(\alph*)]
  \item \emph{(unbiasedness)} Under Assumptions 1.1--1.3,
    $\E(\vec{b} \mid \vec{X}) = \bm{\beta}$.

  \item \emph{(expression for the variance)} Under Assumptions 1.1--1.4,
    $\Var(\vec{b} \mid \vec{X}) = \sigma^{2} \cdot (\vec{X}'\vec{X})^{-1}$.

  \item \emph{(Gauss-Markov Theorem)} Under Assumptions 1.1--1.4, the OLS estimator is
    \textbf{efficient} in the class of linear unbiased estimators.  That is, for any
    unbiased estimator $\widehat{\bm{\beta}}$ that is linear in $\vec{y}$,
    $\Var(\widehat{\bm{\beta}} \mid \vec{X}) \geq \Var(\vec{b} \mid \vec{X})$
    in the matrix sense.\footnote{Let $\vec{A}$ and $\vec{B}$ be two square matrices
    of the same size.  We say that $\vec{A} \geq \vec{B}$ if $\vec{A} - \vec{B}$ is
    positive semidefinite.  A $K \times K$ matrix $\vec{C}$ is said to be positive
    semidefinite (or nonnegative definite) if $\vec{x}'\vec{C}\vec{x} \geq 0$ for all
    $K$-dimensional vectors $\vec{x}$.}

  \item Under Assumptions 1.1--1.4,
    $\Cov(\vec{b}, \vec{e} \mid \vec{X}) = \vec{0}$, where
    $\vec{e} \equiv \vec{y} - \vec{X}\vec{b}$.
\end{enumerate}
\end{proposition}

Before plunging into the proof, let us be clear about what this proposition means.

\begin{itemize}
  \item The matrix inequality in part (c) says that the $K \times K$ matrix
    $\Var(\widehat{\bm{\beta}} \mid \vec{X}) - \Var(\vec{b} \mid \vec{X})$ is positive
    semidefinite, so
    %
    \[
      \vec{a}'[\Var(\widehat{\bm{\beta}} \mid \vec{X})
        - \Var(\vec{b} \mid \vec{X})]\vec{a} \geq 0
      \quad\text{or}\quad
      \vec{a}' \Var(\widehat{\bm{\beta}} \mid \vec{X})\vec{a}
        \geq \vec{a}' \Var(\vec{b} \mid \vec{X})\vec{a}
    \]
    %
    for any $K$-dimensional vector $\vec{a}$.  In particular, consider a special vector
    whose elements are all 0 except for the $k$-th element, which is~1.  For this
    particular $\vec{a}$, the quadratic form $\vec{a}'\vec{A}\vec{a}$ picks up the
    $(k,k)$ element of $\vec{A}$.  But the $(k,k)$ element of
    $\Var(\widehat{\bm{\beta}} \mid \vec{X})$, for example, is
    $\Var(\widehat{\beta}_{k} \mid \vec{X})$ where $\widehat{\beta}_{k}$ is the $k$-th
    element of $\widehat{\bm{\beta}}$.  Thus the matrix inequality in (c) implies
    %
    \begin{equation}
      \Var(\widehat{\beta}_{k} \mid \vec{X})
        \geq \Var(b_{k} \mid \vec{X})
      \quad (k = 1, 2, \ldots, K).
      \label{eq:1.3.1}
    \end{equation}
    %
    That is, for any regression coefficient, the variance of the OLS estimator is no
    larger than that of any other linear unbiased estimator.

  \item As clear from (1.2.5), the OLS estimator is linear in $\vec{y}$.  There are many
    other estimators of $\bm{\beta}$ that are linear and unbiased (you will be asked to
    provide one in a review question below).  The Gauss-Markov Theorem says that the OLS
    estimator is \textbf{efficient} in the sense that its conditional variance matrix
    $\Var(\vec{b} \mid \vec{X})$ is smallest among linear unbiased estimators.  For this
    reason the OLS estimator is called the Best Linear Unbiased Estimator (\textbf{BLUE}).

  \item The OLS estimator $\vec{b}$ is a function of the sample $(\vec{y}, \vec{X})$.
    Since $(\vec{y}, \vec{X})$ are random, so is $\vec{b}$.  Now imagine that we fix
    $\vec{X}$ at some given value, calculate $\vec{b}$ for all samples corresponding to
    all possible realizations of $\vec{y}$, and take the average of $\vec{b}$ (the Monte
    Carlo exercise to this chapter will ask you to do this).  This average is the
    (population) conditional mean $\E(\vec{b} \mid \vec{X})$.  Part (a) (unbiasedness)
    says that this average equals the true value $\bm{\beta}$.

  \item There is another notion of unbiasedness that is weaker than the unbiasedness of
    part (a).  By the Law of Total Expectations,
    $\E[\E(\vec{b} \mid \vec{X})] = \E(\vec{b})$.  So (a) implies
    %
    \begin{equation}
      \E(\vec{b}) = \bm{\beta}.
      \label{eq:1.3.2}
    \end{equation}
    %
    This says: if we calculated $\vec{b}$ for all possible different samples, differing
    not only in $\vec{y}$ but also in $\vec{X}$, the average would be the true value.
    This unconditional statement is probably more relevant in economics because samples do
    differ in both $\vec{y}$ and $\vec{X}$.  The import of the conditional statement (a)
    is that it implies the unconditional statement \eqref{eq:1.3.2}, which is more
    relevant.

  \item The same holds for the conditional statement (c) about the variance.  A review
    question below asks you to show that statements (a) and (b) imply
    %
    \begin{equation}
      \Var(\widehat{\bm{\beta}}) \geq \Var(\vec{b})
      \label{eq:1.3.3}
    \end{equation}
    %
    where $\widehat{\bm{\beta}}$ is any linear unbiased estimator (so that
    $\E(\widehat{\bm{\beta}} \mid \vec{X}) = \bm{\beta}$).
\end{itemize}

We will now go through the proof of this important result.  The proof may look
lengthy; if so, it is only because it records every step, however easy.  In the first
reading, you can skip the proof of part (c).  Proof of (d) is a review question.

\begin{proof}
\begin{enumerate}[label=(\alph*)]

%% ── Part (a): Unbiasedness ──────────────────────────────────────────────────
\item (Proof that $\E(\vec{b} \mid \vec{X}) = \bm{\beta}$)\quad
  $\E(\vec{b} \mid \vec{X}) = \bm{\beta}$ whenever
  $\E(\vec{b} - \bm{\beta} \mid \vec{X}) = \vec{0}$.  So we prove the former.  By the
  expression for the sampling error (1.2.14),
  %
  \[
    \vec{b} - \bm{\beta} = \vec{A}\bm{\varepsilon}
  \]
  %
  where $\vec{A}$ here is $(\vec{X}'\vec{X})^{-1}\vec{X}'$.  So
  %
  \[
    \E(\vec{b} - \bm{\beta} \mid \vec{X})
      = \E(\vec{A}\bm{\varepsilon} \mid \vec{X})
      = \vec{A}\,\E(\bm{\varepsilon} \mid \vec{X}).
  \]
  %
  Here, the second equality holds by the linearity of conditional expectations;
  $\vec{A}$ is a function of $\vec{X}$ and so can be treated as if nonrandom.  Since
  $\E(\bm{\varepsilon} \mid \vec{X}) = \vec{0}$, the last expression is zero.

%% ── Part (b): Expression for the variance ───────────────────────────────────
\item (Proof that $\Var(\vec{b} \mid \vec{X}) = \sigma^{2} \cdot (\vec{X}'\vec{X})^{-1}$)
  %
  \begin{align*}
    \Var(\vec{b} \mid \vec{X})
      &= \Var(\vec{b} - \bm{\beta} \mid \vec{X})
        &\quad&\text{(since $\bm{\beta}$ is not random)} \\
      &= \Var(\vec{A}\bm{\varepsilon} \mid \vec{X})
        &&\text{(by (1.2.14) and $\vec{A} \equiv (\vec{X}'\vec{X})^{-1}\vec{X}'$)} \\
      &= \vec{A}\,\Var(\bm{\varepsilon} \mid \vec{X})\,\vec{A}'
        &&\text{(since $\vec{A}$ is a function of $\vec{X}$)} \\
      &= \vec{A}\,\E(\bm{\varepsilon}\bm{\varepsilon}' \mid \vec{X})\,\vec{A}'
        &&\text{(by Assumption 1.2)} \\
      &= \vec{A}(\sigma^{2}\vec{I}_{n})\vec{A}'
        &&\text{(by Assumption 1.4, see (1.1.14))} \\
      &= \sigma^{2}\vec{A}\vec{A}' \\
      &= \sigma^{2} \cdot (\vec{X}'\vec{X})^{-1}
        &&\text{(since $\vec{A}\vec{A}' = (\vec{X}'\vec{X})^{-1}\vec{X}'\vec{X}(\vec{X}'\vec{X})^{-1} = (\vec{X}'\vec{X})^{-1}$).}
  \end{align*}

%% ── Part (c): Gauss-Markov ──────────────────────────────────────────────────
\item (Gauss-Markov)\quad Since $\widehat{\bm{\beta}}$ is linear in $\vec{y}$, it can
  be written as $\widehat{\bm{\beta}} = \vec{C}\vec{y}$ for some matrix $\vec{C}$,
  which possibly is a function of $\vec{X}$.  Let
  $\vec{D} \equiv \vec{C} - \vec{A}$ or $\vec{C} = \vec{D} + \vec{A}$ where
  $\vec{A} \equiv (\vec{X}'\vec{X})^{-1}\vec{X}'$.  Then
  %
  \begin{align*}
    \widehat{\bm{\beta}}
      &= (\vec{D} + \vec{A})\vec{y} \\
      &= \vec{D}\vec{y} + \vec{A}\vec{y} \\
      &= \vec{D}(\vec{X}\bm{\beta} + \bm{\varepsilon}) + \vec{b}
        &\quad&\text{(since $\vec{y} = \vec{X}\bm{\beta} + \bm{\varepsilon}$ and
        $\vec{A}\vec{y} = (\vec{X}'\vec{X})^{-1}\vec{X}'\vec{y} = \vec{b}$)} \\
      &= \vec{D}\vec{X}\bm{\beta} + \vec{D}\bm{\varepsilon} + \vec{b}.
  \end{align*}
  %
  Taking the conditional expectation of both sides, we obtain
  %
  \[
    \E(\widehat{\bm{\beta}} \mid \vec{X})
      = \vec{D}\vec{X}\bm{\beta}
        + \E(\vec{D}\bm{\varepsilon} \mid \vec{X})
        + \E(\vec{b} \mid \vec{X}).
  \]
  %
  Since both $\vec{b}$ and $\widehat{\bm{\beta}}$ are unbiased and since
  $\E(\vec{D}\bm{\varepsilon} \mid \vec{X})
    = \vec{D}\,\E(\bm{\varepsilon} \mid \vec{X}) = \vec{0}$,
  it follows that $\vec{D}\vec{X}\bm{\beta} = \vec{0}$.  For this to be true for any
  given $\bm{\beta}$, it is necessary that $\vec{D}\vec{X} = \vec{0}$.  So
  $\widehat{\bm{\beta}} = \vec{D}\bm{\varepsilon} + \vec{b}$ and
  %
  \begin{align*}
    \widehat{\bm{\beta}} - \bm{\beta}
      &= \vec{D}\bm{\varepsilon} + (\vec{b} - \bm{\beta}) \\
      &= (\vec{D} + \vec{A})\bm{\varepsilon}
        &\quad&\text{(by (1.2.14)).}
  \end{align*}
  %
  So
  %
  \begin{align*}
    \Var(\widehat{\bm{\beta}} \mid \vec{X})
      &= \Var(\widehat{\bm{\beta}} - \bm{\beta} \mid \vec{X}) \\
      &= \Var[(\vec{D} + \vec{A})\bm{\varepsilon} \mid \vec{X}] \\
      &= (\vec{D} + \vec{A})\,\Var(\bm{\varepsilon} \mid \vec{X})\,
         (\vec{D}' + \vec{A}') \\
      &\qquad\text{(since both $\vec{D}$ and $\vec{A}$ are functions of $\vec{X}$)} \\
      &= \sigma^{2} \cdot (\vec{D} + \vec{A})(\vec{D}' + \vec{A}')
        &\quad&\text{(since $\Var(\bm{\varepsilon} \mid \vec{X}) = \sigma^{2}\vec{I}_{n}$)} \\
      &= \sigma^{2} \cdot
         (\vec{D}\vec{D}' + \vec{A}\vec{D}' + \vec{D}\vec{A}' + \vec{A}\vec{A}').
  \end{align*}
  %
  But $\vec{D}\vec{A}' = \vec{D}\vec{X}(\vec{X}'\vec{X})^{-1} = \vec{0}$ since
  $\vec{D}\vec{X} = \vec{0}$.  Also, $\vec{A}\vec{A}' = (\vec{X}'\vec{X})^{-1}$ as
  shown in (b).  So
  %
  \begin{align*}
    \Var(\widehat{\bm{\beta}} \mid \vec{X})
      &= \sigma^{2} \cdot [\vec{D}\vec{D}' + (\vec{X}'\vec{X})^{-1}] \\
      &\geq \sigma^{2} \cdot (\vec{X}'\vec{X})^{-1}
        &\quad&\text{(since $\vec{D}\vec{D}'$ is positive semidefinite)} \\
      &= \Var(\vec{b} \mid \vec{X})
        &&\text{(by (b)).}
  \end{align*}
\end{enumerate}
\end{proof}

It should be emphasized that the strict exogeneity assumption (Assumption~1.2) is
critical for proving unbiasedness.  Anything short of strict exogeneity will not do.
For example, it is not enough to assume that $\E(\varepsilon_{i} \mid \vec{x}_{i}) = 0$
for all $i$ or that $\E(\vec{x}_{i} \cdot \varepsilon_{i}) = \vec{0}$ for all~$i$.  We
noted in Section~1.1 that most time-series models do not satisfy strict exogeneity even
if they satisfy weaker conditions such as the orthogonality condition
$\E(\vec{x}_{i} \cdot \varepsilon_{i}) = \vec{0}$.  It follows that for those models the
OLS estimator is not unbiased.

%% ─── Finite-Sample Properties of s² ──────────────────────────────────────────
\subsection*{Finite-Sample Properties of \texorpdfstring{$s^{2}$}{s2}}

We defined the OLS estimator of $\sigma^{2}$ in (1.2.13).  It, too, is unbiased.

\begin{proposition}[Unbiasedness of $s^{2}$]
\label{prop:1.2}
\emph{Under Assumptions 1.1--1.4,}
$\E(s^{2} \mid \vec{X}) = \sigma^{2}$
\emph{(and hence $\E(s^{2}) = \sigma^{2}$), provided $n > K$ (so that $s^{2}$ is
well-defined).}
\end{proposition}

We can prove this proposition easily by the use of the trace
operator.\footnote{The \textbf{trace} of a square matrix $\vec{A}$ is the sum of
the diagonal elements of $\vec{A}$:
$\tr(\vec{A}) = \sum_{i} a_{ii}$.}

\begin{proof}
Since $s^{2} = \vec{e}'\vec{e}/(n - K)$, the proof amounts to showing that
$\E(\vec{e}'\vec{e} \mid \vec{X}) = (n - K)\sigma^{2}$.  As shown in (1.2.12),
$\vec{e}'\vec{e} = \bm{\varepsilon}'\vec{M}\bm{\varepsilon}$ where $\vec{M}$ is the
annihilator.  The proof consists of proving two properties:
(1)~$\E(\bm{\varepsilon}'\vec{M}\bm{\varepsilon} \mid \vec{X})
  = \sigma^{2} \cdot \tr(\vec{M})$, and
(2)~$\tr(\vec{M}) = n - K$.

\medskip\noindent
(1) (Proof that
$\E(\bm{\varepsilon}'\vec{M}\bm{\varepsilon} \mid \vec{X})
  = \sigma^{2} \cdot \tr(\vec{M})$)\quad
Since $\bm{\varepsilon}'\vec{M}\bm{\varepsilon}
  = \sum_{i=1}^{n}\sum_{j=1}^{n} m_{ij}\,\varepsilon_{i}\,\varepsilon_{j}$
(this is just writing out the quadratic form $\bm{\varepsilon}'\vec{M}\bm{\varepsilon}$),
we have
%
\begin{align*}
  \E(\bm{\varepsilon}'\vec{M}\bm{\varepsilon} \mid \vec{X})
    &= \sum_{i=1}^{n}\sum_{j=1}^{n} m_{ij}\;\E(\varepsilon_{i}\,\varepsilon_{j}
       \mid \vec{X})
    &\quad&\text{(because $m_{ij}$'s are functions of $\vec{X}$,} \\
    & &&\phantom{\text{(}}\text{$\E(m_{ij}\,\varepsilon_{i}\,\varepsilon_{j} \mid \vec{X})
       = m_{ij}\;\E(\varepsilon_{i}\,\varepsilon_{j} \mid \vec{X})$)} \\
    &= \sum_{i=1}^{n} m_{ii}\,\sigma^{2}
    &&\text{(since $\E(\varepsilon_{i}\,\varepsilon_{j} \mid \vec{X}) = 0$
       for $i \neq j$ by Assumption 1.4)} \\
    &= \sigma^{2} \sum_{i=1}^{n} m_{ii} \\
    &= \sigma^{2} \cdot \tr(\vec{M}).
\end{align*}

\medskip\noindent
(2) (Proof that $\tr(\vec{M}) = n - K$)
%
\begin{align*}
  \tr(\vec{M})
    &= \tr(\vec{I}_{n} - \vec{P})
      &\quad&\text{(since $\vec{M} \equiv \vec{I}_{n} - \vec{P}$; see (1.2.8))} \\
    &= \tr(\vec{I}_{n}) - \tr(\vec{P})
      &&\text{(fact: the trace operator is linear)} \\
    &= n - \tr(\vec{P}),
\end{align*}
%
and
%
\begin{align*}
  \tr(\vec{P})
    &= \tr[\vec{X}(\vec{X}'\vec{X})^{-1}\vec{X}']
      &\quad&\text{(since $\vec{P} \equiv \vec{X}(\vec{X}'\vec{X})^{-1}\vec{X}'$;
      see (1.2.7))} \\
    &= \tr[(\vec{X}'\vec{X})^{-1}\vec{X}'\vec{X}]
      &&\text{(fact: $\tr(\vec{A}\vec{B}) = \tr(\vec{B}\vec{A})$)} \\
    &= \tr(\vec{I}_{K}) = K.
\end{align*}
%
So $\tr(\vec{M}) = n - K$.
\end{proof}

%% ─── Estimate of Var(b | X) ─────────────────────────────────────────────────
\subsection*{Estimate of \texorpdfstring{$\Var(\vec{b} \mid \vec{X})$}{Var(b|X)}}

If $s^{2}$ is the estimate of $\sigma^{2}$, a natural estimate of
$\Var(\vec{b} \mid \vec{X}) = \sigma^{2} \cdot (\vec{X}'\vec{X})^{-1}$ is
%
\begin{equation}
  \widehat{\Var(\vec{b} \mid \vec{X})}
    \equiv s^{2} \cdot (\vec{X}'\vec{X})^{-1}.
  \label{eq:1.3.4}
\end{equation}
%
This is one of the statistics included in the computer printout of any OLS software
package.
